\documentclass[12pt,a4paper]{article}

% ---- Encoding & Font ----
\usepackage{iftex}
\ifPDFTeX
  \usepackage[utf8]{inputenc}
  \usepackage[T1]{fontenc}
\else
  \usepackage{fontspec}  % XeTeX/LuaTeX (tectonic): Unicode nativo (·, —, ¿, ª, §)
\fi
\usepackage[spanish]{babel}
\usepackage{csquotes}

% ---- Page layout ----
\usepackage[top=2.5cm, bottom=2.5cm, left=2.5cm, right=2.5cm]{geometry}
\usepackage{setspace}
\onehalfspacing

% ---- Math ----
\usepackage{amsmath, amssymb, amsthm}
\usepackage{mathtools}
\usepackage{physics}

% ---- Graphics ----
\usepackage{graphicx}
\usepackage{tikz}
\usepackage{float}
\usepackage{caption}

% ---- Tables ----
\usepackage{array}
\usepackage{booktabs}
\usepackage{tabularx}
\usepackage{colortbl}

% ---- Colours ----
\usepackage{xcolor}
\definecolor{notebg}{HTML}{FFF3CD}
\definecolor{noteborder}{HTML}{FFC107}
\definecolor{boxred}{HTML}{F8D7DA}
\definecolor{boxredborder}{HTML}{F5C6CB}

% ---- Boxes ----
\usepackage{mdframed}
\usepackage{tcolorbox}
\tcbuselibrary{most}

\newtcolorbox{keybox}[1][]{
  colback=blue!5,
  colframe=blue!70!black,
  arc=4pt,
  boxrule=1pt,
  fonttitle=\bfseries,
  title={#1}
}

\newtcolorbox{warnbox}[1][]{
  colback=boxred,
  colframe=boxredborder,
  coltitle=black!75,
  arc=4pt,
  boxrule=1pt,
  fonttitle=\bfseries,
  title={#1}
}

\newtcolorbox{notebox}[1][]{
  colback=notebg,
  colframe=noteborder,
  coltitle=black!75,
  arc=4pt,
  boxrule=1pt,
  fonttitle=\bfseries,
  title={#1}
}

% ---- Hyperlinks & TOC ----
\usepackage[hidelinks]{hyperref}
\usepackage{bookmark}

% ---- Enumerate ----
\usepackage{enumitem}

% ---- Miscellaneous ----
\usepackage{icomma}
\usepackage{siunitx}
\usepackage{textcomp}
\usepackage[bottom]{footmisc}

% ---- Diagramas (gráficas) ----
\usepackage{pgfplots}
\pgfplotsset{compat=1.18}
\usetikzlibrary{babel}  % babel español activa `>`; esto evita romper las flechas `->`
% courses/MetNum2023/Clases/assets/tikzstyles.tex
% ---------------------------------------------------------------------------
% Estilos compartidos para los diagramas del curso Métodos Numéricos
% (MetNum2023). Fuente única del "look" visual (colores, grosores, escalas)
% para que todas las figuras (matrices triangulares, curvas de convergencia
% / error, splines a trozos) sean consistentes entre clases.
%
% Uso: la clase debe cargar antes `tikz` y `pgfplots` (bloque §6.3 del
%      CLAUDE.md del curso), y luego
%      % courses/MetNum2023/Clases/assets/tikzstyles.tex
% ---------------------------------------------------------------------------
% Estilos compartidos para los diagramas del curso Métodos Numéricos
% (MetNum2023). Fuente única del "look" visual (colores, grosores, escalas)
% para que todas las figuras (matrices triangulares, curvas de convergencia
% / error, splines a trozos) sean consistentes entre clases.
%
% Uso: la clase debe cargar antes `tikz` y `pgfplots` (bloque §6.3 del
%      CLAUDE.md del curso), y luego
%      % courses/MetNum2023/Clases/assets/tikzstyles.tex
% ---------------------------------------------------------------------------
% Estilos compartidos para los diagramas del curso Métodos Numéricos
% (MetNum2023). Fuente única del "look" visual (colores, grosores, escalas)
% para que todas las figuras (matrices triangulares, curvas de convergencia
% / error, splines a trozos) sean consistentes entre clases.
%
% Uso: la clase debe cargar antes `tikz` y `pgfplots` (bloque §6.3 del
%      CLAUDE.md del curso), y luego
%      \input{../assets/tikzstyles.tex}
% (compilar desde el directorio de la clase para que la ruta relativa resuelva).
%
% Paleta propia de este curso (no se reusa la de Fisica3-2015/CDIV2017/
% ElecMag2024): mismos tonos base para alinear con keybox/notebox/warnbox,
% pero archivo independiente, sin dependencia cruzada entre cursos.
% ---------------------------------------------------------------------------

% ---- Paleta (alineada con las cajas del preámbulo: keybox/notebox/warnbox) ----
\definecolor{figblue}{HTML}{1F4E79}   % azul principal (L, curvas principales, keybox)
\definecolor{figred}{HTML}{B03A2E}    % rojo de énfasis (U, warnbox, error)
\definecolor{figamber}{HTML}{B8860B}  % ámbar (notebox, pivotes)
\definecolor{figgray}{HTML}{555555}   % auxiliares, ejes, ceros de matriz

% ---- TikZ: librerías y estilos reutilizables ----
% Nota: las puntas se declaran como `-{Latex[...]}` y no como `->`. La forma
% con llaves NO contiene un `>` literal, así que es inmune al `>` activo que
% introduce babel español (de todos modos se carga \usetikzlibrary{babel} en
% el bloque §6.3 para cubrir cualquier `->` suelto).
\usetikzlibrary{arrows.meta,calc,matrix,patterns}

\tikzset{
  figline/.style={figblue, line width=0.9pt},
  figaux/.style={figgray, dashed, line width=0.6pt},
  figaxis/.style={figgray, line width=0.7pt, -{Latex[length=2.2mm,width=1.5mm]}},
  % estructura de matrices (L triangular inferior, U triangular superior)
  matcelda/.style={minimum width=7mm, minimum height=7mm, anchor=center, font=\small},
  matnz/.style={matcelda, fill=figblue!12},
  matnzU/.style={matcelda, fill=figred!12},
  matcero/.style={matcelda, text=figgray!70},
  matpivote/.style={matcelda, fill=figamber!25, draw=figamber, line width=0.8pt},
  % nodos de interpolación / splines
  fignodo/.style={circle, inner sep=0pt, minimum size=4.5pt, draw=figblue, line width=0.9pt, fill=white},
  fignodoactivo/.style={fignodo, fill=figamber},
  % etiquetas
  figlbl/.style={font=\small},
  figlblaux/.style={font=\footnotesize, figgray},
}

% ---- pgfplots: estilos base ----
\pgfplotsset{
  % convergencia / error vs. n (grado, iteración) — suele ir en escala log
  errorfig/.style={
    width=10cm, height=6.5cm,
    xlabel style={font=\small},
    ylabel style={font=\small},
    tick label style={font=\footnotesize},
    every axis plot/.append style={line width=1.1pt, mark=*, mark size=1.5pt},
    grid=major, grid style={figgray!15},
    legend cell align=left,
    legend style={font=\footnotesize, draw=figgray!40, at={(0.5,-0.25)}, anchor=north},
    clip=false,
  },
  % interpolación / splines: un color por tramo, nodos marcados aparte
  interpfig/.style={
    width=10cm, height=6cm,
    axis lines=middle,
    xlabel style={font=\small, at={(axis description cs:1,0.5)}, anchor=west},
    ylabel style={font=\small, at={(axis description cs:0.5,1)}, anchor=south},
    tick label style={font=\footnotesize},
    every axis plot/.append style={line width=1.1pt},
    grid=none,
    legend cell align=left,
    legend style={font=\footnotesize, draw=figgray!40},
    clip=false,
  },
}

% (compilar desde el directorio de la clase para que la ruta relativa resuelva).
%
% Paleta propia de este curso (no se reusa la de Fisica3-2015/CDIV2017/
% ElecMag2024): mismos tonos base para alinear con keybox/notebox/warnbox,
% pero archivo independiente, sin dependencia cruzada entre cursos.
% ---------------------------------------------------------------------------

% ---- Paleta (alineada con las cajas del preámbulo: keybox/notebox/warnbox) ----
\definecolor{figblue}{HTML}{1F4E79}   % azul principal (L, curvas principales, keybox)
\definecolor{figred}{HTML}{B03A2E}    % rojo de énfasis (U, warnbox, error)
\definecolor{figamber}{HTML}{B8860B}  % ámbar (notebox, pivotes)
\definecolor{figgray}{HTML}{555555}   % auxiliares, ejes, ceros de matriz

% ---- TikZ: librerías y estilos reutilizables ----
% Nota: las puntas se declaran como `-{Latex[...]}` y no como `->`. La forma
% con llaves NO contiene un `>` literal, así que es inmune al `>` activo que
% introduce babel español (de todos modos se carga \usetikzlibrary{babel} en
% el bloque §6.3 para cubrir cualquier `->` suelto).
\usetikzlibrary{arrows.meta,calc,matrix,patterns}

\tikzset{
  figline/.style={figblue, line width=0.9pt},
  figaux/.style={figgray, dashed, line width=0.6pt},
  figaxis/.style={figgray, line width=0.7pt, -{Latex[length=2.2mm,width=1.5mm]}},
  % estructura de matrices (L triangular inferior, U triangular superior)
  matcelda/.style={minimum width=7mm, minimum height=7mm, anchor=center, font=\small},
  matnz/.style={matcelda, fill=figblue!12},
  matnzU/.style={matcelda, fill=figred!12},
  matcero/.style={matcelda, text=figgray!70},
  matpivote/.style={matcelda, fill=figamber!25, draw=figamber, line width=0.8pt},
  % nodos de interpolación / splines
  fignodo/.style={circle, inner sep=0pt, minimum size=4.5pt, draw=figblue, line width=0.9pt, fill=white},
  fignodoactivo/.style={fignodo, fill=figamber},
  % etiquetas
  figlbl/.style={font=\small},
  figlblaux/.style={font=\footnotesize, figgray},
}

% ---- pgfplots: estilos base ----
\pgfplotsset{
  % convergencia / error vs. n (grado, iteración) — suele ir en escala log
  errorfig/.style={
    width=10cm, height=6.5cm,
    xlabel style={font=\small},
    ylabel style={font=\small},
    tick label style={font=\footnotesize},
    every axis plot/.append style={line width=1.1pt, mark=*, mark size=1.5pt},
    grid=major, grid style={figgray!15},
    legend cell align=left,
    legend style={font=\footnotesize, draw=figgray!40, at={(0.5,-0.25)}, anchor=north},
    clip=false,
  },
  % interpolación / splines: un color por tramo, nodos marcados aparte
  interpfig/.style={
    width=10cm, height=6cm,
    axis lines=middle,
    xlabel style={font=\small, at={(axis description cs:1,0.5)}, anchor=west},
    ylabel style={font=\small, at={(axis description cs:0.5,1)}, anchor=south},
    tick label style={font=\footnotesize},
    every axis plot/.append style={line width=1.1pt},
    grid=none,
    legend cell align=left,
    legend style={font=\footnotesize, draw=figgray!40},
    clip=false,
  },
}

% (compilar desde el directorio de la clase para que la ruta relativa resuelva).
%
% Paleta propia de este curso (no se reusa la de Fisica3-2015/CDIV2017/
% ElecMag2024): mismos tonos base para alinear con keybox/notebox/warnbox,
% pero archivo independiente, sin dependencia cruzada entre cursos.
% ---------------------------------------------------------------------------

% ---- Paleta (alineada con las cajas del preámbulo: keybox/notebox/warnbox) ----
\definecolor{figblue}{HTML}{1F4E79}   % azul principal (L, curvas principales, keybox)
\definecolor{figred}{HTML}{B03A2E}    % rojo de énfasis (U, warnbox, error)
\definecolor{figamber}{HTML}{B8860B}  % ámbar (notebox, pivotes)
\definecolor{figgray}{HTML}{555555}   % auxiliares, ejes, ceros de matriz

% ---- TikZ: librerías y estilos reutilizables ----
% Nota: las puntas se declaran como `-{Latex[...]}` y no como `->`. La forma
% con llaves NO contiene un `>` literal, así que es inmune al `>` activo que
% introduce babel español (de todos modos se carga \usetikzlibrary{babel} en
% el bloque §6.3 para cubrir cualquier `->` suelto).
\usetikzlibrary{arrows.meta,calc,matrix,patterns}

\tikzset{
  figline/.style={figblue, line width=0.9pt},
  figaux/.style={figgray, dashed, line width=0.6pt},
  figaxis/.style={figgray, line width=0.7pt, -{Latex[length=2.2mm,width=1.5mm]}},
  % estructura de matrices (L triangular inferior, U triangular superior)
  matcelda/.style={minimum width=7mm, minimum height=7mm, anchor=center, font=\small},
  matnz/.style={matcelda, fill=figblue!12},
  matnzU/.style={matcelda, fill=figred!12},
  matcero/.style={matcelda, text=figgray!70},
  matpivote/.style={matcelda, fill=figamber!25, draw=figamber, line width=0.8pt},
  % nodos de interpolación / splines
  fignodo/.style={circle, inner sep=0pt, minimum size=4.5pt, draw=figblue, line width=0.9pt, fill=white},
  fignodoactivo/.style={fignodo, fill=figamber},
  % etiquetas
  figlbl/.style={font=\small},
  figlblaux/.style={font=\footnotesize, figgray},
}

% ---- pgfplots: estilos base ----
\pgfplotsset{
  % convergencia / error vs. n (grado, iteración) — suele ir en escala log
  errorfig/.style={
    width=10cm, height=6.5cm,
    xlabel style={font=\small},
    ylabel style={font=\small},
    tick label style={font=\footnotesize},
    every axis plot/.append style={line width=1.1pt, mark=*, mark size=1.5pt},
    grid=major, grid style={figgray!15},
    legend cell align=left,
    legend style={font=\footnotesize, draw=figgray!40, at={(0.5,-0.25)}, anchor=north},
    clip=false,
  },
  % interpolación / splines: un color por tramo, nodos marcados aparte
  interpfig/.style={
    width=10cm, height=6cm,
    axis lines=middle,
    xlabel style={font=\small, at={(axis description cs:1,0.5)}, anchor=west},
    ylabel style={font=\small, at={(axis description cs:0.5,1)}, anchor=south},
    tick label style={font=\footnotesize},
    every axis plot/.append style={line width=1.1pt},
    grid=none,
    legend cell align=left,
    legend style={font=\footnotesize, draw=figgray!40},
    clip=false,
  },
}


% ---- Title ----
\title{\textbf{Clase 13: Interpolación a Trozos} \\[0.3em]
        \large{Interpolación lineal y cúbica a trozos, fenómeno de Runge, interpolación de Hermite}}
\author{Transcripto y expandido --- Curso Métodos Numéricos (OpenFING)}
\date{}

\begin{document}

\maketitle
\thispagestyle{empty}
\newpage

\tableofcontents
\newpage

\section{Motivación: el fenómeno de Runge}

La clase retoma el teorema de error de interpolación probado la clase
anterior (Clase 12): dado un polinomio interpolante $p_n$ de grado $\le n$
por $n+1$ nodos $x_0,\dots,x_n \in [a,b]$,
\[
f(x) - p_n(x) = \frac{f^{(n+1)}(\gamma_x)}{(n+1)!}\,\omega_n(x), \qquad \omega_n(x) = \prod_{i=0}^n (x-x_i),
\]
para cierto $\gamma_x$ no controlable. La cota resultante depende de tres
factores: la derivada de orden $n+1$ de $f$, el factorial $(n+1)!$ (que
ayuda, porque crece muy rápido) y el \textbf{polinomio nodal} $\omega_n$
(que puede ayudar o no, según cómo se elijan los nodos).

Se retoma el ejemplo de la \textbf{función de Runge},
\[
f(x) = \frac{1}{1+25x^2}, \qquad x \in [-1,1],
\]
interpolada con nodos equiespaciados de grado creciente. A diferencia del
ejemplo anterior ($\mathrm{sen}(x)\cos(x)$, donde la interpolante convergía
muy bien), acá la interpolante \textbf{no converge} en norma del supremo:
para $n$ grande aparecen oscilaciones violentas cerca de los extremos del
intervalo (el fenómeno se ve con claridad al graficar, aunque en el centro
del intervalo la aproximación es buena).

\begin{notebox}[Por qué pasa esto]
Con nodos equiespaciados, el polinomio nodal $\omega_n$ no ayuda ni
perjudica ``a propósito'': simplemente puede volverse enorme. Para la
función de Runge, las derivadas de orden alto crecen tan rápido que ni el
factorial $(n+1)!$ del denominador alcanza para compensar, y la cota (y el
error real) diverge.
\end{notebox}

\noindent\textbf{Mensaje general:} si una función tiene derivadas de orden
alto que crecen rápido, hacer interpolación polinomial de grado alto es mala
idea.

\section{Dos salidas: mejores nodos o interpolar a trozos}

Como el factor $f^{(n+1)}$ y el factorial $(n+1)!$ no son negociables (son
datos del problema), la única palanca disponible es el polinomio nodal
$\omega_n$. Se mencionan dos estrategias:

\begin{enumerate}
  \item \textbf{Elegir mejor los nodos} --- los \textbf{nodos de Chebyshev}
  (construidos proyectando una partición uniforme de un círculo sobre el
  diámetro) minimizan $\|\omega_n\|_\infty$ y evitan el fenómeno de Runge
  cerca de los bordes. La clase los muestra brevemente (con 5, 20 y 50 nodos
  la curva mejora notoriamente), pero \textbf{no se desarrollan más}: quedan
  mencionados como una alternativa válida, sin demostración ni fórmula de
  error propia en esta clase.
  \item \textbf{Interpolar a trozos con grado bajo} --- en vez de un único
  polinomio de grado alto (con mucha ``libertad para oscilar''), particionar
  el intervalo en subintervalos chicos y usar un polinomio de grado bajo en
  cada uno. Caso extremo: una función \textbf{lineal} no puede tener un
  máximo local, así que no puede oscilar. Ésta es la estrategia que
  desarrolla el resto de la clase.
\end{enumerate}

\begin{notebox}[Cuándo vale la pena]
Para una función ``linda'' como $\mathrm{sen}(x)\cos(x)$, interpolar a
trozos no aporta nada --- conviene el interpolante global de grado alto. La
interpolación a trozos se justifica cuando \textbf{no se pueden elegir los
nodos} (o cuando elegirlos como Chebyshev no alcanza) y la función tiene
derivadas de orden alto grandes.
\end{notebox}

\section{Interpolación lineal a trozos}

\subsection{Definición}

Dados $(x_i, y_i)$, $i=0,\dots,n$, con $x_0 < x_1 < \dots < x_n$, se define
la interpolante lineal a trozos $L: [x_0,x_n] \to \mathbb{R}$ como la
función que, restringida a cada subintervalo $[x_i, x_{i+1}]$, es el (único)
polinomio de grado $\le 1$ que interpola $(x_i,y_i)$ y $(x_{i+1},y_{i+1})$:
\[
L(x) = y_i + (x - x_i)\,\frac{y_{i+1}-y_i}{x_{i+1}-x_i}, \qquad x \in [x_i, x_{i+1}].
\]
Es literalmente ``unir los puntos con segmentos de recta'' --- la
existencia no plantea ningún problema porque ya se sabe interpolar por dos
puntos.

\subsection{Continuidad sin derivabilidad}

$L$ es \textbf{continua} por construcción (los segmentos comparten
extremos), pero \textbf{no es derivable} en los nodos interiores en general
(salvo que $f$ sea constante) --- hay un ``pico'' en cada $x_i$.

\subsection{Cota del error}

Si los $(x_i,y_i)$ provienen de evaluar una función $f$, se reutiliza el
teorema de error de interpolación de la clase anterior, pero
\textbf{aplicado localmente}: para $x \in [x_i,x_{i+1}]$ existe $\gamma_x
\in (x_i,x_{i+1})$ tal que
\[
f(x) - L(x) = \frac{f''(\gamma_x)}{2}\,(x-x_i)(x-x_{i+1}).
\]
Es exactamente el teorema de la clase anterior particularizado a $n=1$ en el
subintervalo $[x_i,x_{i+1}]$ en vez de en todo $[a,b]$.

Acotando: $(x-x_i)(x-x_{i+1})$ (con signo cambiado para que sea positivo) es
una parábola que vale $0$ en ambos extremos y, por simetría, alcanza su
máximo en el punto medio, donde vale $\frac{(x_{i+1}-x_i)^2}{4}$. Tomando el
peor caso sobre todos los subintervalos:
\begin{keybox}[Cota del error --- interpolación lineal a trozos]
\[
\|f - L\|_\infty \;\le\; \frac{\|f''\|_{\infty,[x_0,x_n]}}{8}\,h^2,
\qquad h = \max_i (x_{i+1}-x_i).
\]
\end{keybox}

\noindent\textit{De dónde sale el 8:} sale de combinar el $2$ del $2!$ en el
denominador del teorema con el $4$ del máximo de la parábola
$(x-x_i)(x-x_{i+1})$.

\subsection{Aplicación a la función de Runge}

Con $n+1$ nodos equiespaciados en $[-1,1]$, $h = \dfrac{2}{n}$. La derivada
segunda de la función de Runge,
\[
f''(x) = \frac{50\,(75x^2-1)}{(1+25x^2)^3},
\]
alcanza su máximo en valor absoluto en $x=0$, donde vale $50$ (dato que se
da sin demostrarlo formalmente en la clase --- se justifica
cualitativamente mirando el gráfico: la derivada segunda mide ``cuán rápido
cambia la pendiente'', y eso es máximo en el centro, donde la función pasa
de crecer a decrecer más bruscamente; hacia los extremos la función se
parece a una recta y $f''$ es chica).
\[
\|f - L\|_\infty \;\le\; \frac{50}{8}\cdot\frac{4}{n^2} \;=\; \frac{25}{n^2}
\;\xrightarrow[n\to\infty]{}\; 0.
\]

% assets/clase13-runge-lineal.tex
% Función de Runge f(x)=1/(1+25x^2) en [-1,1] vs. su interpolante lineal a
% trozos con 9 nodos equiespaciados. Ilustra §3.4 del summary: a diferencia
% de la interpolación global de grado alto (Clase 12), la lineal a trozos
% no oscila cerca de los bordes.
\begin{figure}[H]
  \centering
  \begin{tikzpicture}
    \begin{axis}[interpfig,
      xlabel={$x$}, ylabel={$y$},
      xmin=-1.05, xmax=1.05, ymin=-0.05, ymax=1.12,
      % el eje x va abajo (no centrado): rótulos en la punta de cada eje
      xlabel style={at={(ticklabel* cs:1)}, anchor=west},
      ylabel style={at={(ticklabel* cs:1)}, anchor=south},
      % leyenda debajo del eje, fuera de la curva
      legend style={at={(0.5,-0.2)}, anchor=north}, legend columns=2,
      height=5.2cm,
    ]
      \addplot[figblue, line width=1.1pt, domain=-1:1, samples=200, forget plot]
        {1/(1+25*x^2)};
      \addlegendimage{figblue, line width=1.1pt}
      \addlegendentry{$f(x)=1/(1+25x^2)$}

      \addplot[figred, line width=1pt, mark=none]
        coordinates {
          (-1,0.038462) (-0.75,0.066390) (-0.5,0.137931) (-0.25,0.390244)
          (0,1) (0.25,0.390244) (0.5,0.137931) (0.75,0.066390) (1,0.038462)
        };
      \addlegendentry{lineal a trozos ($n=8$)}

      \addplot[only marks, mark=*, mark size=2pt, mark options={fill=figamber, draw=figblue}, forget plot]
        coordinates {
          (-1,0.038462) (-0.75,0.066390) (-0.5,0.137931) (-0.25,0.390244)
          (0,1) (0.25,0.390244) (0.5,0.137931) (0.75,0.066390) (1,0.038462)
        };
    \end{axis}
  \end{tikzpicture}
  \caption{Función de Runge y su interpolante lineal a trozos con 9 nodos equiespaciados. A diferencia del interpolante polinomial global (Clase 12), no hay oscilación cerca de los bordes.}
  \label{fig:clase13-runge-lineal}
\end{figure}


\begin{keybox}[Contraste clave con la clase anterior]
Con los mismos nodos equiespaciados que producían el fenómeno de Runge en la
interpolación global de grado alto, la interpolación \textbf{lineal a
trozos} converge uniformemente. El precio es una curva visualmente ``tosca''
(llena de picos), pero cumple el objetivo de convergencia en norma infinito.
\end{keybox}

\section{De lineal a cúbica a trozos}

\subsection{Por qué no cuadrática}

Se busca ahora un interpolante a trozos \textbf{derivable}, no solo
continuo: que en cada nodo interior $x_i$ coincidan tanto el valor como la
derivada por izquierda y por derecha. Para un intervalo interior
$[x_i,x_{i+1}]$ eso son \textbf{cuatro condiciones}:
\[
p(x_i)=y_i,\quad p(x_{i+1})=y_{i+1},\quad p'(x_i)=d_i,\quad p'(x_{i+1})=d_{i+1}.
\]
Un polinomio cuadrático tiene sólo \textbf{tres} grados de libertad --- no
alcanzan para las cuatro condiciones. Un polinomio \textbf{cúbico} tiene
exactamente cuatro coeficientes, así que el sistema queda determinado (y es
a lo sumo un sistema $4\times 4$, resoluble por ejemplo a la Vandermonde,
aunque no se resuelve en el pizarrón).

\subsection{Por qué grado impar}

En general estos interpolantes a trozos se toman de \textbf{grado impar}:
un polinomio de grado impar tiene una cantidad \textbf{par} de coeficientes,
lo que encaja con imponer pares de condiciones (valor + derivadas de
distintos órdenes) en cada extremo del subintervalo.

\begin{notebox}
Se menciona además, sin desarrollarlo, que una ventaja práctica de trabajar
a trozos (frente a un interpolante global) es que agregar un nodo nuevo
sólo obliga a recalcular el (o los) subintervalo(s) afectados, no toda la
interpolante.
\end{notebox}

\section{La interpolante cúbica a trozos}

\subsection{El problema, con los \texorpdfstring{$d_i$}{di} dados}

Se plantea el problema suponiendo, por ahora, que además de los datos
$(x_i,y_i)$ también se conocen valores deseados de la derivada en cada
nodo, $d_i$ (de dónde salen esos $d_i$ es la pregunta que se responde en la
sección \ref{sec:tres-formas}). El problema es: hallar
$P:[x_0,x_n]\to\mathbb{R}$ tal que, en cada subintervalo, $P$ sea un
polinomio de grado $\le 3$ con
\[
P(x_i)=y_i, \qquad P'(x_i)=d_i \qquad \text{para todo } i.
\]

\subsection{Fórmula explícita (base de Hermite)}

Con $s = x-x_i$ la variable local dentro del subintervalo y $h_i =
x_{i+1}-x_i$ su longitud, la interpolante en $[x_i,x_{i+1}]$ es
\begin{keybox}[Interpolante cúbica a trozos]
\[
P(x) = y_i\,\frac{(h_i-s)^2(h_i+2s)}{h_i^3}
     \;+\; y_{i+1}\,\frac{s^2(3h_i-2s)}{h_i^3}
     \;+\; d_i\,\frac{s(s-h_i)^2}{h_i^2}
     \;+\; d_{i+1}\,\frac{s^2(s-h_i)}{h_i^2}.
\]
\end{keybox}

\noindent\textit{Verificación en los extremos} (como se hace en clase): en
$s=0$ sobreviven sólo los términos sin $s$, y queda $P=y_i$; en $s=h_i$
sobrevive sólo el término de $y_{i+1}$, y queda $P=y_{i+1}$ --- consistente
con lo pedido. (Puede verificarse igual de mecánicamente que
$P'(x_i)=d_i$ y $P'(x_{i+1})=d_{i+1}$, aunque la clase no hace esa cuenta
explícita.)

Es, en palabras del docente, ``una de esas fórmulas que escribís una vez,
la copiás y la pegás'': una vez fijados los $d_i$ (por cualquiera de los
tres métodos de la sección \ref{sec:tres-formas}), esta fórmula da el
interpolante cúbico a trozos completo.

\section{Tres formas de elegir los \texorpdfstring{$d_i$}{di}}
\label{sec:tres-formas}

La clase anuncia (y desarrolla en las próximas clases, ``entre hoy y el
miércoles'') tres estrategias distintas para fijar los $d_i$, con preguntas
distintas asociadas a cada una:

\begin{enumerate}
  \item \textbf{Interpolación de Hermite} --- los $d_i$ \emph{vienen
  dados}: si los datos provienen de una función conocida $f$, se toma
  $d_i = f'(x_i)$. La pregunta relevante es cuán bien aproxima el
  interpolante a $f$ (se desarrolla en la sección \ref{sec:hermite}).
  \item \textbf{Splines} --- los $d_i$ \textbf{no} vienen dados; se eligen
  de forma que la función resultante tenga \textbf{derivada segunda
  continua} (sea de clase $C^2$). No se desarrolla en esta clase.
  \item \textbf{Interpolación que preserva forma (\texttt{pchip})} --- los
  datos son puntos sueltos, sin función subyacente conocida; los $d_i$ se
  eligen con fórmulas heurísticas para que la curva ``se vea linda'' (es lo
  que hace la función \texttt{pchip} de MATLAB/Octave). No se desarrolla en
  esta clase; para estos dos últimos casos no se plantea una pregunta de
  error contra una $f$ (no hay $f$), sino de estética de la curva.
\end{enumerate}

\section{Interpolación de Hermite cúbica a trozos}
\label{sec:hermite}

\subsection{Definición}

Dada $f:[x_0,x_n]\to\mathbb{R}$ derivable y nodos $x_0<\dots<x_n$, la
\textbf{interpolante de Hermite cúbica a trozos} es la función $P$ dada en
cada subintervalo por la fórmula de la sección anterior con $y_i = f(x_i)$
y $d_i = f'(x_i)$. Por construcción, $P$ es continua, con derivada
continua, coincide con $f$ en los nodos y su derivada coincide con $f'$ en
los nodos.

\subsection{Teorema del error}

A diferencia de la interpolante lineal a trozos, acá se dispone de más
información (los valores de la derivada), y la fórmula de error de la
Clase 12 ``se queda corta'': sólo explota la coincidencia de valores, no de
derivadas. El teorema (enunciado primero para un único intervalo $[a,b]$,
antes de pegar subintervalos):

\begin{keybox}[Teorema --- error de la interpolación de Hermite cúbica]
Sea $f$ de clase $C^4$ en $[a,b]$ y sea $P$ el polinomio de grado $\le 3$
tal que $P(a)=f(a)$, $P(b)=f(b)$, $P'(a)=f'(a)$, $P'(b)=f'(b)$. Entonces
para todo $x\in[a,b]$ existe $\gamma_x \in (a,b)$ tal que
\[
f(x)-P(x) \;=\; \frac{f^{(4)}(\gamma_x)}{24}\,(x-a)^2(x-b)^2.
\]
\end{keybox}

Comparado con el teorema de la clase anterior (con nodos simples), acá el
polinomio nodal aparece \textbf{al cuadrado}, aparece la derivada
\textbf{cuarta} (no la $n+1$-ésima genérica) y el denominador es $24$ (no
un factorial genérico). El docente es explícito en que el valor concreto
del ``24'' no importa memorizarlo --- lo que importa es la técnica de la
demostración.

\subsection{Idea de la demostración}

Es la misma técnica que en la Clase 12 (función auxiliar + Rolle
repetido), con un ajuste. Se define, para el $x$ fijado,
\[
g(t) = f(t) - P(t) - E(x)\cdot \frac{(t-a)^2(t-b)^2}{(x-a)^2(x-b)^2}, \qquad E(x) := f(x)-P(x),
\]
donde el factor $(t-a)^2(t-b)^2$ (en vez de $(t-a)(t-b)$, como en la clase
anterior) es la clave del ajuste: al elevarlo al cuadrado, \textbf{su
derivada también se anula} en $a$ y en $b$ (no sólo la función).
Concretamente:

\begin{itemize}
  \item $g$ se anula en $a$, $b$ y en $x$ $\to$ \textbf{tres} raíces de
  $g$.
  \item Además $g'(a)=0$ y $g'(b)=0$: el término $f'-P'$ se cancela porque
  $P'(a)=f'(a)$ y $P'(b)=f'(b)$ (por definición de $P$), y el término que
  viene del factor cuadrático también se anula en $a$ y $b$ porque su
  derivada tiene un factor $(t-a)(t-b)$.
\end{itemize}

Con esto, $g'$ tiene \textbf{cuatro} raíces en $(a,b)$ (dos entre $a$-$x$ y
$x$-$b$ por Rolle aplicado a $g$, más $a$ y $b$ mismos como raíces directas
de $g'$). Aplicando Rolle sucesivamente: $g''$ tiene al menos 3 raíces,
$g'''$ al menos 2, y $g^{(4)}$ al menos 1 --- llamada $\gamma_x$. Igualando
$g^{(4)}(\gamma_x)=0$ y despejando se obtiene la fórmula del teorema.

\begin{notebox}[La diferencia con la demostración anterior]
Está toda concentrada en ese cambio del factor lineal al cuadrático: ganar
dos raíces adicionales de $g'$ (en $a$ y en $b$) es lo que permite llegar
hasta la derivada cuarta en vez de quedarse en la $(n+1)$-ésima genérica
del caso simple.
\end{notebox}

\vspace{1em}
\noindent\textit{Clase que viene: splines cúbicos (elegir los $d_i$ para
$C^2$) e interpolación que preserva forma (\texttt{pchip}).}

\end{document}
